% Created 2026-11-18 Wed 11:40
% Intended LaTeX compiler: xelatex
\documentclass[12pt]{article}
\usepackage{graphicx}
\usepackage{longtable}
\usepackage{wrapfig}
\usepackage{rotating}
\usepackage[normalem]{ulem}
\usepackage{amsmath}
\usepackage{amssymb}
\usepackage{capt-of}
\usepackage{hyperref}
\usepackage[a4paper,margin=1in]{geometry}
\usepackage{xurl}
\author{fcalado}
\date{\today}
\title{}
\hypersetup{
 pdfauthor={fcalado},
 pdftitle={},
 pdfkeywords={},
 pdfsubject={},
 pdfcreator={Emacs 29.3 (Org mode 9.6.15)}, 
 pdflang={English}}
\begin{document}


\section*{INFO 656: Machine Learning}
\label{sec:org507f431}

\section*{Final Project}
\label{sec:org6d4bcaa}

\subsection*{Instructions}
\label{sec:org7f011cb}

For the final project, you will do an evaluation of a machine learning
dataset, model, or both. For example, you may prompt the same model
over two or more different datasets, or prompt two different models
with the same (or multiple) datasets, in order to evaluate changes in
the outputs.

You may evaluate any aspect of the model's outputs, like social bias,
performance, accuracy, or another metric.

Regardless of the model and dataset that you choose to evaluate, you
must include a close reading analysis of model outputs. Close reading
means that you select individual examples to compare and examine the
outputs closely, paying attention to word choice, tone, repetition,
connotation, etc. You must choose at least 4 examples for close
reading analysis, and each example should include a couple or more
sentences of analysis.

Please refer to the in-class example of evaluation, which you may use
as a template for your chosen dataset and/or models:
\url{https://github.com/gofilipa/656/blob/main/final-project-eval-template.ipynb}

In your close reading analysis, you might draw from useful examples
from our course readings, like Devinney, Hannah, et al.
\url{https://doi.org/10.18653/v1/2024.gebnlp-1.3} (for a useful example of
what close reading of model outputs might look like); Abeba Birhane et
al, "The Dark Side of Dataset Scaling: Evaluating Racial
Classification in Multimodal Models."
\url{https://doi.org/10.1145/3630106.3658968} (for a critical perspective of
training datasets).

If you choose, you may also do a quantitative evaluation of the
outputs, though this is not required. See the Huggingface Evaluate
Library for more information about quantitative evaluation:
\url{https://huggingface.co/docs/evaluate/choosing\_a\_metric}

\subsection*{Deliverables}
\label{sec:orgb84a760}

By the deadline of \textbf{\textbf{Wednesday, December 9}} at 3pm, You will be
required to submit the following:
\begin{itemize}
\item a python notebook performing the evaluation, on Canvas
\item a close reading analysis of at least 4 outputs (2 for each dataset,
if you are comparing 2 datasets, for example), submitted as a
separate document, on Canvas
\item an overall assessment or main finding(s) from your close reading
analysis (submitted along with the close reading document on Canvas)
\item a short, 5 minute presentation on your process and main finding(s),
in person
\end{itemize}
\end{document}
